\documentclass[10pt,a4paper]{article}

\usepackage[utf8]{inputenc}
\usepackage[T1]{fontenc}
\usepackage[english]{babel}
\usepackage{amsmath,amssymb}
\usepackage[top=1.1cm,bottom=1.1cm,left=1.7cm,right=1.7cm]{geometry}
\usepackage{titlesec}
\usepackage{hyperref}
\hypersetup{colorlinks=true,linkcolor=black,citecolor=blue,urlcolor=blue}

\setlength{\parskip}{0.15em}
\titlespacing*{\section}{0pt}{0.7em}{0.3em}

\title{\vspace{-1.2cm}Bibliographic Review, Topic 39\\[0.2em]
\large The Blind Spot Paradox}
\author{Alexandre Boyer \and Mélaine Gouillou \and Salomé Fonvielle \and Ulysse Petit-Tichanné\\[0.2em]
\normalsize Research Track, CentraleSupélec. Supervisor: Raphaël Minato}
\date{}

\begin{document}
\maketitle
\vspace{-1.4em}

\section{The problem setting}

A model deployed on a stream must cope with \emph{concept drift}, a change over time in
the relation $P(y\mid X)$. Moreno-Torres et al.~\cite{moreno2012} separate it from a shift
in $P(X)$ alone, in a taxonomy they build because the literature had been using the same
words for different things. The two surveys that frame the field, Gama et al.~\cite{gama2014}
and Lu et al.~\cite{lu2018}, describe the same deployed architecture: an adaptive classifier
paired with an external detector watching its error rate. Both treat that pair as
cooperative, and neither asks whether the classifier might destroy what the detector is
trying to measure.

Each block has its own reaction time. The Adaptive Random Forest~\cite{gomes2017} is a bag
of Hoeffding Trees~\cite{domingos2000}, incremental and inert once converged, resampled by
Poisson$(\lambda{=}6)$ weights, and every tree carries its own ADWIN. ADWIN~\cite{bifet2007}
splits a variable-length window in two and declares a change when the sub-window means
diverge past a Hoeffding bound; on a warning the ARF grows a background tree, and on a
confirmation it swaps that tree in. The external block is a CUSUM~\cite{page1954} whose
Page-Hinkley form accumulates $S_t=\max(0,\,S_{t-1}+e_t-p_0-\delta_P)$ and fires past
$\lambda$. Hinkley~\cite{hinkley1971} supplies the companion result, the asymptotic law of
the change point \emph{estimated} from that same statistic. By Wald's equation the
accumulation time is $\tau^{*}_{\det}=\lambda/(\Delta e-\delta_P)=O(1/\Delta e)$.

Both blocks read the same binary error stream, and one erases what the other needs to see.
The question matters in finance, where regime switches are abrupt and volatility clusters.
ProteuS~\cite{suarez2025} fits ARMA--GARCH models to four ETFs and interpolates between
them, producing synthetic returns with a known break date, the ground truth that real
market data cannot supply.

\section{What the source paper establishes}

Working in River 0.23.0~\cite{montiel2021} and instrumenting the ARF's internal swap counter
rather than simulating the race with proxy streams, the manuscript establishes three things.

\paragraph{A clock artefact, not an ADWIN flaw.} River evaluates ADWIN every $c$ steps,
$c=32$ by default, while the ARF's internal detectors run at $c=1$. The measured delay is
then $\mathrm{ADD}=c-1-(\tau^{*}\bmod c)+\tau_{\mathrm{stat}}$, uniform on
$\{\tau_{\mathrm{stat}},\dots,\tau_{\mathrm{stat}}+c-1\}$ for an arbitrary drift date. The
default pipeline handicaps the external monitor before any statistics are involved.

\paragraph{Starvation, amplified by the Hydra effect.} The forest adapts as soon as its
fastest tree does, $\tau_{\mathrm{ARF}}=\min_{i\le M}\tau_i$, and the instrumentation
measures that as $4$ to $8$ times faster than a single tree
($\tau_{\mathrm{ARF}}\approx18.5\,\Delta e^{-0.98}$ against
$\tau_{\mathrm{HAT}}\approx102\,\Delta e^{-1.02}$: same exponent, smaller prefactor). The
error excess therefore lasts $O(1/\Delta e)$, and so does the CUSUM's accumulation
requirement. The two scale together, so a violent shock is repaired faster and leaves the
detector no more evidence than a slow one; a Doob maximal inequality over the transient
window bounds the crossing probability away from~$1$. On ProteuS streams, a PHT paired with
an ARF at $c_{\mathrm{int}}{=}1$ detects $0$ of $1{,}080$ runs, whereas throttling the
forest to $c_{\mathrm{int}}{=}32$ restores $F_1\approx0.6$. The same collapse appears with
EDDM and with Streaming Random Patches, which is what makes it a property of the
cumulative-evidence family rather than of one detector.

\paragraph{Consequences.} Under an assumption of independent trees,
$P_{\mathrm{miss}}=1-\bigl(1-F(\tau^{*}_{\det})\bigr)^{M}$, whence a critical ensemble size
$M_{\mathrm{crit}}=\lfloor\ln\beta/\ln(1-F(\tau^{*}_{\det}))\rfloor$ beyond which no
guarantee survives. At $\Delta e=0.33$ it evaluates to $0$, so the certificate fails even
for a single tree. The proposed fix, the Decoupling Principle, requires
$c_{\mathrm{ext}}<c_{\mathrm{int}}$ and $\lambda<\lambda^{*}(\Delta e_{\min})$, and it runs
straight into a negative result the authors state plainly: on heteroscedastic streams,
decoupling caps $\lambda$ at $12.4$ while false-alarm control demands at least $15$, which
leaves the admissible set empty. Detection is recovered by changing detector family rather
than by tuning. KSWIN~\cite{raab2020}, a Kolmogorov--Smirnov test between sliding windows,
keeps pre-drift data in its reference window and cannot be starved.

\section{Where our reading departs from the manuscript}

\paragraph{The race is between two random times.} The manuscript compares
$\tau_{\mathrm{ARF}}$ against the nominal $\tau^{*}_{\det}=\lambda/(\Delta e-\delta_P)$, a
deterministic quantity, whereas the alarm is itself a stopping time that may never fire.
The natural object is the cumulative incidence function of competing risks,
$F_1(t)=\Pr(T\le t,\ \varepsilon=\text{repair})$ with
$T=\min(\tau_{\mathrm{ARF}},\tau_{\det})$. Fine and Gray~\cite{fine1999} model its
subdistribution hazard directly, their argument being that the cause-specific hazard
carries no usable interpretation in terms of the probability of failing from that cause.
Estimation goes through the competing-risks case of the Aalen--Johansen product-limit
estimator~\cite{aalen1978}, not a Kaplan--Meier treating alarms as independent censoring:
an alarm structurally prevents observing the repair in the same run, which is the
dependence Kaplan--Meier assumes away.

\paragraph{The independence assumption is optimistic, and the bias can be signed.} The
$\tau_i$ enter through an order statistic~\cite{david2003}, but the $M$ trees share one
stream and are not independent. Conditioning on the stream realization $S$ leaves only the
per-tree Poisson weights and feature subspaces, which are i.i.d.\ across trees. Writing
$G_S(s)=P(\tau_1>s\mid S)$ and applying Jensen's inequality to the convex map
$x\mapsto x^{M}$ gives $P(\tau_{\mathrm{ARF}}>s)=\mathbb{E}_S[G_S(s)^M]\ge(1-F(s))^{M}$.
The real repair is slower than the formula predicts, with equality only if the difficulty
of adapting does not depend on which stream was drawn, a case that does not arise here
since every seed produces a different post-drift noise trajectory. Esary, Proschan and
Walkup~\cite{esary1967} reach the same place from the other side: their Theorem~5.1 gives
$\Pr(\tau_1>s_1,\dots,\tau_M>s_M)\ge\prod_i\Pr(\tau_i>s_i)$ for associated variables, and
the law of total covariance shows that the trees here can only be non-negatively
associated. The manuscript concedes the point in a disclaimer; making it quantitative turns
$M_{\mathrm{crit}}$ into a conservative floor rather than an estimate.

\paragraph{No standard metric exists for what has to be measured.} Both surveys enumerate
the evaluation toolkit, prequential error, kappa and kappa-temporal statistics, prequential
AUC, RAM-hours~\cite{gama2014,lu2018}, and neither offers a time-to-recovery measure. That
gap is why the group instruments the swap fraction, the smoothed ensemble error and the
detector's internal state together, and why $\tau_{\mathrm{ARF}}$, the date of the first
swap alone, has to be validated as a proxy for ensemble recovery instead of assumed to be
one.

\vspace{-0.2em}
{\footnotesize
\begin{thebibliography}{99}\setlength{\itemsep}{0pt}
\bibitem{moreno2012} J.~G. Moreno-Torres, T.~Raeder, R.~Alaiz-Rodríguez, N.~V. Chawla, and F.~Herrera, ``A unifying view on dataset shift in classification,'' \emph{Pattern Recognition}, vol.~45, no.~1, pp.~521--530, 2012.
\bibitem{gama2014} J.~Gama, I.~Žliobaitė, A.~Bifet, M.~Pechenizkiy, and A.~Bouchachia, ``A survey on concept drift adaptation,'' \emph{ACM Computing Surveys}, vol.~46, no.~4, art.~44, 2014.
\bibitem{lu2018} J.~Lu, A.~Liu, F.~Dong, F.~Gu, J.~Gama, and G.~Zhang, ``Learning under concept drift: a review,'' \emph{IEEE Trans. Knowl. Data Eng.}, vol.~31, no.~12, pp.~2346--2363, 2018.
\bibitem{gomes2017} H.~M. Gomes, A.~Bifet, J.~Read, J.~P. Barddal, F.~Enembreck, B.~Pfahringer, G.~Holmes, and T.~Abdessalem, ``Adaptive random forests for evolving data stream classification,'' \emph{Machine Learning}, vol.~106, no.~9--10, pp.~1469--1495, 2017.
\bibitem{domingos2000} P.~Domingos and G.~Hulten, ``Mining high-speed data streams,'' in \emph{Proc. KDD}, 2000, pp.~71--80.
\bibitem{bifet2007} A.~Bifet and R.~Gavaldà, ``Learning from time-changing data with adaptive windowing,'' in \emph{Proc. SIAM SDM}, 2007, pp.~443--448.
\bibitem{page1954} E.~S. Page, ``Continuous inspection schemes,'' \emph{Biometrika}, vol.~41, no.~1--2, pp.~100--115, 1954.
\bibitem{hinkley1971} D.~V. Hinkley, ``Inference about the change-point from cumulative sum tests,'' \emph{Biometrika}, vol.~58, no.~3, pp.~509--523, 1971.
\bibitem{suarez2025} A.~L. Suárez-Cetrulo, A.~Cervantes, and D.~Quintana, ``ProteuS: a generative approach for simulating concept drift in financial markets,'' arXiv:2509.11844, 2025.
\bibitem{montiel2021} J.~Montiel, M.~Halford, S.~M. Mastelini et al., ``River: machine learning for streaming data in Python,'' \emph{JMLR}, vol.~22, pp.~1--8, 2021.
\bibitem{raab2020} C.~Raab, M.~Heusinger, and F.-M. Schleif, ``Reactive soft prototype computing for concept drift streams,'' \emph{Neurocomputing}, vol.~416, pp.~340--351, 2020.
\bibitem{fine1999} J.~P. Fine and R.~J. Gray, ``A proportional hazards model for the subdistribution of a competing risk,'' \emph{JASA}, vol.~94, no.~446, pp.~496--509, 1999.
\bibitem{aalen1978} O.~O. Aalen and S.~Johansen, ``An empirical transition matrix for non-homogeneous Markov chains based on censored observations,'' \emph{Scand. J. Statist.}, vol.~5, no.~3, pp.~141--150, 1978.
\bibitem{david2003} H.~A. David and H.~N. Nagaraja, \emph{Order Statistics}, 3rd ed. Wiley, 2003. (Not available online; the association argument of~\cite{esary1967} was used instead.)
\bibitem{esary1967} J.~D. Esary, F.~Proschan, and D.~J. Walkup, ``Association of random variables, with applications,'' \emph{Ann. Math. Statist.}, vol.~38, no.~5, pp.~1466--1474, 1967.
\end{thebibliography}
}

\end{document}
